\documentclass[english, aspectratio=169]{beamer}
% english is for the language used in standard texts (figures, tables etc)
% aspectratio of 16:9 or set it for more old school to 4:3 (without the ':')

% ---------------------------------------------------------------------------- %
% Load base preamble
% ---------------------------------------------------------------------------- %
\usepackage{import}
\subimport{./preamble/}{beamer.tex}

% ---------------------------------------------------------------------------- %
% Local settings
% ---------------------------------------------------------------------------- %
\newcommand{\B}[0]{\ensuremath{\mathbb{B}}}

\newcommand{\sort}[0]{\text{sort}}

\newcommand{\triple}[3]{\ensuremath{(#1, #2, #3)}}
\renewcommand{\arc}[3]{\ensuremath{#1 \xrightarrow{_{#2}} #3}}


\tikzstyle{plot_adiar}=[color=black, mark=o, mark size=1pt, line width=0.7pt]
\tikzstyle{plot_buddy}=[color=red, mark=o, mark size=1pt, line width=0.7pt]
\tikzstyle{plot_cudd}=[color=blue, mark=diamond, mark size=1pt, line width=0.7pt]
\tikzstyle{plot_sylvan}=[color=purple, mark=square, mark size=1pt, line width=0.7pt]

% Horizontal legends: https://tex.stackexchange.com/a/101578
% argument #1: any options
\makeatletter
\newenvironment{customlegend}[1][]{%
  \begingroup
  % inits/clears the lists (which might be populated from previous
  % axes):
  \pgfplots@init@cleared@structures
  \pgfplotsset{#1}%
}{%
  % draws the legend:
  \pgfplots@createlegend
  \endgroup
}%

% makes \addlegendimage available (typically only available within an
% axis environment):
\def\addlegendimage{\pgfplots@addlegendimage}
\makeatother

% ------------------------------------------------------------------------------
% 
% ------------------------------------------------------------------------------
% Opener: ...
% 
% Key points:
% Existential
% Basics (single variable)
% Push
% Bounce
% Relabelling
% Bounce
% 
% Take Home Message: How to do I/O-efficient "multi-recursion" (for RelProd)

% ------------------------------------------------------------------------------
% TITLEPAGE
% ------------------------------------------------------------------------------
\title{An External Memory Relational Product}

\author{\textbf{Steffan Christ S\o lvsten}, Jaco van de Pol}

\institute{\includegraphics[width=0.2\linewidth]{external/aulogo_uk_var2_black.eps}}

\date{\today}

\begin{document}
\titleframe

\begin{frame}
  \begin{center}
    \LARGE

    $\mathit{RelProd}(S,T)
    \equiv
    (
    \ {\only<2>{\color{orange}} \exists \vec{x} .}\
    S(\vec{x}) \wedge T(\vec{x}, \vec{x'})
    \ ) [\vec{x'} / \vec{x}]$
  \end{center}
\end{frame}

\begin{frame}
  \frametitle{$\exists x_i .\ \phi \equiv \phi[x_i / \top] \vee \phi[x_i / \bot]$}

  \begin{center}
  \begin{tikzpicture}
    % Sweepline
    \onslide<2> { \draw[gray, thick, dashed] (-3, 0) -- (10, 0); }
    \onslide<3> { \draw[gray, thick, dashed] (-3, -0.9) -- (10, -0.9); }
    \onslide<4> { \draw[gray, thick, dashed] (-3, -1.85) -- (10, -1.85); }
    \onslide<5> { \draw[gray, thick, dashed] (-3, -2.9) -- (10, -2.9); }
    \onslide<6> { \draw[gray, thick, dashed] (-3, -3.9) -- (10, -3.9); }
    \onslide<7> { \draw[gray, thick, dashed] (-3, -4.75) -- (10, -4.75); }

    % INPUT
    % TODO: Use shape from (push)?
    \draw[black, thick] (0.1,0.1)  -- (1.8,-1.85)  -- (1.5,-3.9)  -- (0.81,-5.7)
                        (-0.1,0.1) -- (-1.8,-1.8) -- (-1.5,-3.9) -- (-0.8,-5.7);

    \node[shape = circle, draw = black, fill = white]
    (i_root) {};

    \onslide<2-> {
      \node[shape = circle, draw = black, fill = white,
            below right=0.65cm and 0.1cm of i_root]
      (i_p1) {};

      \draw[->] (i_root) edge (i_p1);
    }

    \onslide<3-> {
      \node[shape = circle, draw = black, fill = white,
            below left=0.65cm and 0.1cm of i_p1]
      (i_p2) {};

      \draw[->, dashed] (i_p1) edge (i_p2);
    }

    \onslide<4-> {
      \node[shape = circle, draw = black, fill = white,
            below right=0.7cm and 0.1cm of i_p2]
      (i_quant) {\tiny $x_i$};

      \draw[->] (i_p2) edge (i_quant);
    }

    \onslide<5-> {
      \node[shape = circle, draw = black, fill = white,
            below left =0.65cm and 0.1cm of i_quant]
      (i_p3_1) {};

      \node[shape = circle, draw = black, fill = white,
            below right =0.65cm and 0.1cm of i_quant]
      (i_p3_2) {};

      \draw[->, dashed] (i_quant) edge (i_p3_1);

      \draw[->] (i_quant) edge (i_p3_2);
    }

    \onslide<6-> {
      \node[shape = circle, draw = black, fill = white,
            below left =0.6cm and 0.2cm of i_p3_1]
      (i_p4_1) {};

      \draw[->, dashed] (i_p3_1) edge (i_p4_1);

      \node[shape = circle, draw = black, fill = white,
            below left =0.6cm and 0.2cm of i_p3_2]
      (i_p4_2) {};

      \draw[->, dashed] (i_p3_2) edge (i_p4_2);
    }

    \node[shape = rectangle, draw = black, fill = white,
    below left=5.4cm and 0.2cm of i_root]
    (i_p5_1) {};

    \node[shape = rectangle, draw = black, fill = white,
    below right=5.4cm and 0.2cm of i_root]
    (i_p5_2) {};

    \onslide<7-> {
      \draw[->] (i_p4_1) edge (i_p5_1);
      \draw[->] (i_p4_2) edge (i_p5_2);
    }

    % OUTPUT
    \onslide<3-> {
      \draw[blue, thick] (6.5,0.1)  -- (7.4,-0.9)
                         (6.25,0.1) -- (5.35,-0.9);

      \node[shape = circle, draw = black, fill = white,
            right=6cm of i_root]
      (o_root) {};

      \node[shape = circle, draw = black, fill = white,
            below right=0.65cm and 0.1cm of o_root]
      (o_p1) {};

      \draw[->, blue] (o_p1) edge (o_root);
    }

    \onslide<4-> {
      \draw[blue, thick] (7.4,-0.9)  -- (8.2,-1.85)
                         (5.35,-0.9) -- (4.55,-1.85);

      \node[shape = circle, draw = black, fill = white,
            below left=0.65cm and 0.1cm of o_p1]
      (o_p2) {};

      \draw[->, dashed, blue] (o_p2) edge (o_p1);
    }

    \onslide<5> {
      \draw[blue, dotted, thick] (8.2,-1.85) -- (8.45, -2.9)
                                 (4.55,-1.85) -- (4.27,-2.9);
    }

    \onslide<6-> {
      \draw[blue, thick] (8.2,-1.85) -- (8.7, -3.9)
                         (4.55,-1.85) -- (4,-3.9);

      \node[shape = circle, draw = black, fill = white,
            below right=1.78cm and 0.1cm of o_p2]
       (o_p3) {};

       \draw[->, blue] (o_p3) edge (o_p2);
    }

    \onslide<7-8> {
      \draw[blue, thick] (8.7, -3.9) -- (8.6, -4.75)
                         (4,-3.9) -- (4.2,-4.75);

      \node[shape = circle, draw = black, fill = white,
            below left =0.6cm and 0.1cm of o_p3]
      (o_p4) {};

      \draw[->, dashed, blue] (o_p4) edge (o_p3);

      \node[shape = rectangle, draw = red, fill = white,
            below =5.4cm of o_root]
      (o_p5) {};

      \draw[->, red] (o_p4) edge (o_p5);

      \draw[red, thick] (8.6, -4.75) -- (8,-5.7)
      (4.2,-4.75)  -- (4.8,-5.7);
    }
  \end{tikzpicture}
\end{center}
\end{frame}

\blankframe

\begin{frame}
  \frametitle{$\exists \vec{x}.\ \phi(\dots)$ \quad (Push)}

  \begin{center}
  \begin{tikzpicture}
    % Sweepline
    \onslide<2> { \draw[gray, thick, dashed] (-3, 0) -- (10, 0); }
    \onslide<3> { \draw[gray, thick, dashed] (-3, -0.9) -- (10, -0.9); }
    \onslide<4> { \draw[gray, thick, dashed] (-3, -1.7) -- (10, -1.7); }
    \onslide<5> { \draw[gray, thick, dashed] (-3, -2.5) -- (10, -2.5); }
    \onslide<6> { \draw[gray, thick, dashed] (-3, -3.3) -- (10, -3.3); }
    \onslide<7> { \draw[gray, thick, dashed] (-3, -4.2) -- (10, -4.2); }
    \onslide<8> { \draw[gray, thick, dashed] (-3, -5.0) -- (10, -5.0); }

    % INPUT
    \draw[black, thick] (0.1,0.1)  -- (2,-1.7)  -- (1.8,-4.2)  -- (0.81,-5.7)
                        (-0.1,0.1) -- (-2,-1.7) -- (-1.8,-4.2) -- (-0.8,-5.7);

    \node[shape = circle, draw = black, fill = white]
    (i_root) {};

    \onslide<2->{
      \node[shape = circle, draw = black, fill = white,
            below right=0.55cm and 0.1cm of i_root]
      (i_quant1) {\tiny $x_i$};

      \draw[->] (i_root) edge (i_quant1);
    }

    \onslide<3->{
      \node[shape = circle, draw = black, fill = white,
            below left=0.5cm and 0.2cm of i_quant1]
      (i_p1_1) {};

      \draw[->, dashed] (i_quant1) edge (i_p1_1);

      \node[shape = circle, draw = black, fill = white,
            below right=0.5cm and 0.2cm of i_quant1]
      (i_p1_2) {};

      \draw[->] (i_quant1) edge (i_p1_2);
    }

    \onslide<4->{
      \node[shape = circle, draw = black, fill = white,
            below left=0.5cm and 0.3cm of i_p1_1]
      (i_p2_1) {};

      \node[shape = circle, draw = black, fill = white,
            below left=0.5cm and 0.2cm of i_p1_2]
      (i_p2_2) {};

      \draw[->, dashed]
        (i_p1_1) edge (i_p2_1)
        (i_p1_2) edge (i_p2_2)
      ;
    }

    \onslide<5->{
      \node[shape = circle, draw = black, fill = white,
            below right=0.5cm and 0.1cm of i_p2_1]
      (i_quant2_1) {\tiny $x_j$};

      \node[shape = circle, draw = black, fill = white,
            below right=0.5cm and 0.1cm of i_p2_2]
      (i_quant2_2) {\tiny $x_j$};

      \draw[->]
        (i_p2_1) edge (i_quant2_1)
        (i_p2_2) edge (i_quant2_2)
      ;
    }

    \onslide<6->{
      \node[shape = circle, draw = black, fill = white,
            below left=0.5cm and -0.1cm of i_quant2_1]
      (i_p3_1_1) {};
      \node[shape = circle, draw = black, fill = white,
            below right=0.5cm and -0.1cm of i_quant2_1]
      (i_p3_1_2) {};

      \node[shape = circle, draw = black, fill = white,
            below left=0.5cm and -0.1cm of i_quant2_2]
      (i_p3_2_1) {};
      \node[shape = circle, draw = black, fill = white,
            below right=0.5cm and -0.1cm of i_quant2_2]
      (i_p3_2_2) {};

      \draw[->, dashed]
        (i_quant2_1) edge (i_p3_1_1)
        (i_quant2_2) edge (i_p3_2_1)
      ;

      \draw[->]
        (i_quant2_1) edge (i_p3_1_2)
        (i_quant2_2) edge (i_p3_2_2)
      ;
    }

    \onslide<7->{
      \node[shape = circle, draw = black, fill = white,
            below left=0.5cm and 0.1cm of i_p3_1_1]
      (i_p4_1_1) {};
      \node[shape = circle, draw = black, fill = white,
            below left=0.5cm and 0.1cm of i_p3_1_2]
      (i_p4_1_2) {};

      \node[shape = circle, draw = black, fill = white,
            below left=0.5cm and 0.2cm of i_p3_2_1]
      (i_p4_2_1) {};
      \node[shape = circle, draw = black, fill = white,
            below left=0.5cm and 0.2cm of i_p3_2_2]
      (i_p4_2_2) {};

      \draw[->, dashed]
        (i_p3_1_1) edge (i_p4_1_1)
        (i_p3_1_2) edge (i_p4_1_2)
        (i_p3_2_1) edge (i_p4_2_1)
        (i_p3_2_2) edge (i_p4_2_2)
      ;
    }

    \node[shape = rectangle, draw = black, fill = white,
          below left=5.4cm and 0.1cm of i_root]
    (i_p5_1) {};

    \node[shape = rectangle, draw = black, fill = white,
          below right=5.4cm and 0.1cm of i_root]
    (i_p5_2) {};

    \onslide<8-> {
      \draw[->]
        (i_p4_1_1) edge (i_p5_1)
        (i_p4_1_2) edge (i_p5_1)
        (i_p4_2_1) edge (i_p5_2)
        (i_p4_2_2) edge (i_p5_1)
      ;
    }

    % OUTPUT
    \onslide<3> {
      \draw[blue, dotted, thick] (5.3,-0.9) -- (6.25,0.1) -- (6.45,0.1) -- (7.5,-0.9);
    }

    \onslide<4-> {
      \draw[blue, thick] (6.25,0.1) -- (4.5,-1.7)
                         (6.45,0.1) -- (8.3,-1.7);

      \node[shape = circle, draw = black, fill = white,
        right=6cm of i_root]
      (o_root) {};

      \node[shape = circle, draw = black, fill = white,
            below right=1.4cm and 0.2cm of o_root]
      (o_p1) {};

      \draw[->, blue] (o_p1) edge (o_root);
    }

    \onslide<5-> {
      \draw[blue, thick] (4.5,-1.7) -- (4.3,-2.5)
                         (8.3,-1.7) -- (8.6,-2.5);

      \node[shape = circle, draw = black, fill = white,
            below left=0.55cm and 0.2cm of o_p1]
      (o_p2) {};

      \draw[->, dashed, blue] (o_p2) edge (o_p1);
    }

    \onslide<6-> {
      \draw[blue, thick] (4.3,-2.5) -- (4.35, -3.3)
                         (8.6,-2.5) -- (8.55, -3.3);

      \node[shape = circle, draw = black, fill = white,
        below right=0.5cm and 0.2cm of o_p2]
      (o_quant2) {\tiny $\vee$};

      \draw[->, blue] (o_quant2) edge (o_p2);
    }

    \onslide<7-> {
      \draw[blue, thick] (4.35, -3.3) -- (3.7, -4.2)
                         (8.55, -3.3) -- (9.4, -4.2);

      \node[shape = circle, draw = black, fill = white,
            below left=0.55cm and 0.2cm of o_quant2]
      (o_p3_1) {};

      \draw[->, blue, dashed] (o_p3_1) edge (o_quant2);

      \node[shape = circle, draw = black, fill = white,
            below right=0.55cm and 0.2cm of o_quant2]
      (o_p3_2) {};

      \draw[->, blue] (o_p3_2) edge (o_quant2);
    }

    \onslide<8-9> {
      \draw[blue, thick] (3.7, -4.2) -- (3.9, -5.0)
                         (9.4, -4.2) -- (9.2, -5.0);

      \node[shape = circle, draw = black, fill = white,
            below left=0.55cm and 0.2cm of o_p3_1]
      (o_p4_1) {};

      \node[shape = circle, draw = black, fill = white,
            below left=0.55cm and 0.2cm of o_p3_2]
      (o_p4_2) {};

      \draw[->, blue, dashed]
        (o_p4_1) edge (o_p3_1)
        (o_p4_2) edge (o_p3_2)
      ;

      \draw[red, thick] (3.9, -5.0) -- (4.3, -5.7)
                        (9.2, -5.0) -- (8.7, -5.7);

      \node[shape = rectangle, draw = red, fill = white,
            below right=5.4cm and -0.3cm of o_root]
      (o_p5_1) {};

      \node[shape = rectangle, draw = red, fill = white,
            below right=5.4cm and 0.7cm of o_root]
      (o_p5_2) {};

      \draw[->, red]
        (o_p4_1) edge (o_p5_1)
        (o_p4_2) edge (o_p5_2)
      ;
    }
  \end{tikzpicture}
\end{center}
\end{frame}

\blankframe

\begin{frame}
  \frametitle{$\exists \vec{x}.\ \phi(\dots)$ \quad (Bounce)}

  \begin{center}
  \begin{tikzpicture}
    % Sweepline
    \onslide<2> { \draw[gray, thick, dashed] (-2.5, -5.0) -- (10, -5.0); }
    \onslide<3> { \draw[gray, thick, dashed] (-2.5, -4.2) -- (10, -4.2); }
    \onslide<4-9> { \draw[gray, thick, dashed] (-2.5, -3.2) -- (10, -3.2); }

    \onslide<4>   { \draw[gray, thick, dashed] (2.5, -3.4) -- (9.7, -3.4); }
    \onslide<5,9> { \draw[gray, thick, dashed] (2.5, -4.2) -- (9.7, -4.2); }
    \onslide<6-8> { \draw[gray, thick, dashed] (2.5, -5.0) -- (9.7, -5.0); }

    \onslide<10> { \draw[gray, thick, dashed] (-2.5, -2.5) -- (10, -2.5); }
    \onslide<11> { \draw[gray, thick, dashed] (-2.5, -1.7) -- (10, -1.7); }
    \onslide<12> { \draw[gray, thick, dashed] (-2.5, 0) -- (10, 0); }

    % --------------------------------------------------------------------------
    % INPUT
    \draw[blue, thick] (-1.8,-5.0) -- (-2.0,-4.2) -- (-1.4,-1.7) -- (-0.1,0.1)
                    -- (0.1,0.1)  -- (1.4,-1.7)  -- (2.0,-4.2)  -- (1.8,-5.0);

    \draw[red, thick] (-1.8,-5.0) -- (-1.6,-5.6)
                      (1.8,-5.0) -- (1.6,-5.6);

    \onslide<11-> {
      \node[shape = circle, draw = black, fill = white]
      (i_root) {};
    }

    \onslide<10-> {
      \node[shape = circle, draw = black, fill = white,
            below right=1.4cm and 0.2cm of i_root]
      (i_p1) {};
    }

    \onslide<11-> {
      \draw[->, blue] (i_p1) edge (i_root);
    }

    \onslide<3-> {
      \node[shape = circle, draw = black, fill = white,
            below left=0.55cm and 0.2cm of i_p1]
      (i_p2) {};
    }

    \onslide<10-> {
      \draw[->, dashed, blue] (i_p2) edge (i_p1);
    }

    \onslide<3-> {
      \node[shape = circle, draw = black, fill = white,
            below right=0.5cm and 0.2cm of i_p2]
      (i_quant2) {\tiny $x_j$};
    }

    \onslide<4-> {
      \draw[->, blue, thick] (i_quant2) edge (i_p2);
    }

    \onslide<2-> {
      \node[shape = circle, draw = black, fill = white,
            below left=0.55cm and 0.2cm of i_quant2]
      (i_p3_1) {};

      \node[shape = circle, draw = black, fill = white,
            below right=0.55cm and 0.2cm of i_quant2]
      (i_p3_2) {};
    }

    \onslide<3-> {
      \draw[->, blue, thick, dashed] (i_p3_1) edge (i_p2);
      \draw[->, blue, dashed] (i_p3_1) edge (i_quant2);
      \draw[->, blue] (i_p3_2) edge (i_quant2);
    }

    \onslide<2-> {
      \node[shape = circle, draw = black, fill = white,
            below left=0.55cm and 0.2cm of i_p3_1]
      (i_p4_1) {};

      \node[shape = circle, draw = black, fill = white,
            below left=0.55cm and 0.2cm of i_p3_2]
      (i_p4_2) {};

      \draw[->, blue, dashed]
        (i_p4_1) edge (i_p3_1)
        (i_p4_2) edge (i_p3_2)
      ;
    }

    \node[shape = rectangle, draw = red, fill = white,
      below right=5.4cm and -0.3cm of i_root]
    (i_p5_1) {};

    \node[shape = rectangle, draw = red, fill = white,
      below right=5.4cm and 0.7cm of i_root]
    (i_p5_2) {};

    \onslide<2-> {
      \draw[->, red]
        (i_p4_1) edge (i_p5_1)
        (i_p4_2) edge (i_p5_2)
      ;
    }

    % --------------------------------------------------------------------------
    % REDUCED OUTPUT
    \onslide<2-6,8-> {
      \draw[thick] (3.0,-5.0) -- (3.2,-5.6)
                   (5.3,-5.0) -- (5.1,-5.6)
      ;
          
      \node[shape = rectangle, draw = black, fill = white,
            right=3.5cm of i_p5_1]
      (o1_p5_1) {};

      \node[shape = rectangle, draw = black, fill = white,
            right=3.5cm of i_p5_2]
      (o1_p5_2) {};

      \node[shape = circle, draw = black, fill = white,
            right=3.5cm of i_p4_1]
      (o1_p4_1) {};

      \node[shape = circle, draw = black, fill = white,
            right=3.5cm of i_p4_2]
      (o1_p4_2) {};

      \draw[->]
        (o1_p4_1) edge (o1_p5_1)
        (o1_p4_2) edge (o1_p5_2)
      ;
    }

    \onslide<3-5,9-> {
      \draw[thick] (2.9,-4.2) -- (3.0,-5.0)
                   (5.5,-4.2) -- (5.3,-5.0)
      ;

      \node[shape = circle, draw = black, fill = white,
            right=3.5cm of i_p3_1]
      (o1_p3_1) {};

      \node[shape = circle, draw = black, fill = white,
            right=3.5cm of i_p3_2]
      (o1_p3_2) {};

      \draw[->, dashed]
        (o1_p3_1) edge (o1_p4_1)
        (o1_p3_2) edge (o1_p4_2)
      ;
    }

    \onslide<4> {
      \draw[thick, dotted] (3.2,-3.4) -- (2.9,-4.2)
                           (5.1,-3.4) -- (5.5,-4.2)
      ;
    }

    \onslide<10-> {
      \draw[thick] (2.8,-2.5) -- (2.9,-4.2)
                   (4.8,-2.5) -- (5.5,-4.2)
      ;

      \node[shape = circle, draw = black, fill = white,
            right=3.5cm of i_p2]
      (o1_p2) {};

      \draw[->, dashed] (o1_p2) edge (o1_p3_1);
      \draw[->] (o1_p2) edge (o1_p3_2);
    }

    \onslide<11-> {
      \draw[thick] (2.9,-1.7) -- (2.8,-2.5)
                   (4.6,-1.7) -- (4.8,-2.5)
      ;

      \node[shape = circle, draw = black, fill = white,
            right=3.5cm of i_p1]
      (o1_p1) {};

      \draw[->, dashed] (o1_p1) edge (o1_p2);
    }

    \onslide<12-13> {
      \draw[thick] (3.8,0) -- (2.9,-1.7)
                   (4.0,0) -- (4.6,-1.7)
      ;

      \node[shape = circle, draw = black, fill = white,
            right=3.5cm of i_root]
      (o1_root) {};

      \draw[->] (o1_root) edge (o1_p1);
    }

    % --------------------------------------------------------------------------
    % SECOND SWEEP
    \onslide<5-9> {
      \draw[blue, thick, dotted] (6.95, -2.9) -- (6.9,-3.2)
                                 (9.15, -2.9) -- (9.2,-3.2)
          ;

      \draw[blue, thick] (6.9,-3.2) -- (6.8, -4.2)
                         (9.2,-3.2) -- (9.4, -4.2)
      ;
          
      \node[shape = circle, draw = none, % dummy for crossing in-going arcs
            right=3.5cm of o1_p2]
      (o2_p2) {};

      \node[shape = circle, draw = black, fill = white,
            right=3.5cm of o1_p3_1]
      (o2_p3_1) {};

      \node[shape = circle, draw = black, fill = white,
            right=3.5cm of o1_p3_2]
      (o2_p3_2) {};

      \draw[->, thick, blue, dashed] (o2_p3_1) edge (o2_p2);
      \draw[->, thick, blue] (o2_p3_2) edge (o2_p2);
    }

    \onslide<6-8> {
      \draw[blue, thick] (6.8, -4.2) -- (6.9, -5.0)
                         (9.4, -4.2) -- (9.3, -5.0)
      ;

      \node[shape = circle, draw = black, fill = white,
            right=3.5cm of o1_p4_1]
      (o2_p4_1) {};

      \node[shape = circle, draw = black, fill = white,
            right=3.5cm of o1_p4_2]
      (o2_p4_2) {};

      \draw[->, blue, dashed]
        (o2_p4_1) edge (o2_p3_1)
        (o2_p4_2) edge (o2_p3_2)
      ;

      \draw[red, thick] (6.9, -5.0) -- (7.2, -5.6)
                        (9.3, -5.0) -- (9.0, -5.6)
      ;

      \node[shape = rectangle, draw = red, fill = white,
            right=3.5cm of o1_p5_1]
      (o2_p5_1) {};

      \node[shape = rectangle, draw = red, fill = white,
            right=3.5cm of o1_p5_2]
      (o2_p5_2) {};

      \draw[->, red]
        (o2_p4_1) edge (o2_p5_1)
        (o2_p4_2) edge (o2_p5_2)
      ;
    }
  \end{tikzpicture}
\end{center}
\end{frame}

\blankframe

\begin{frame}
  \begin{center}
    \LARGE

    $\mathit{RelProd}(S,T)
    \equiv
    {\color{orange} (}
    \ \exists \vec{x} .\
    S(\vec{x}) \wedge T(\vec{x}, \vec{x'})
    \ {\color{orange} ) [\vec{x'} / \vec{x}]}$
  \end{center}
\end{frame}

\begin{frame}
  \begin{definition}
    A relabelling $\pi$ is monotonic if $x_i < x_j \implies \pi(x_i) < \pi(x_j)$
  \end{definition}

  \pause
  {\bf If $\pi$ is monotonic}

  \begin{itemize}
  \item \emph{1-Var} / \emph{Push}:

    Apply $\pi$ in $O(L_N)$ extra time during the final bottom-up Reduce
    sweep.

  \item \emph{Bounce}:

    $x_i' < x_j' \implies x_i < x_j$: $\pi$ can be applied during the
    outermost Reduce sweep.
  \end{itemize}

  \quad That is, applying $\pi$ can (essentially) be done for free.

  \pause
  {\bf If $\pi$ is not monotonic}

  \quad to be continued...
\end{frame}
\end{document}
